% !TeX program = xelatex
% !TeX encoding = UTF-8
% ========================================================================
\documentclass{resume}
% ========================================================================
% ========================================================================
\usepackage{zh_CN_fonts}
\usepackage{linespacing_fix}
\usepackage{amsmath,bm}
\usepackage{booktabs}
\usepackage{cite}
\usepackage[hidelinks,bookmarks=false]{hyperref}
\usepackage{fontawesome} % 图标
\usepackage{datetime2} % 日期时间
\usepackage{enumitem} % 用于自定义列表
% 在导言区添加（\begin{document}之前）
\newcommand{\confbadge}[2]{\hfill\makebox[3.5em][r]{\textcolor{#1}{\textbf{[#2]}}}}
% ========================================================================
% ========================================================================
\begin{document}
% ========================================================================
% ========================================================================
\pagenumbering{gobble}
% ========================================================================
\begin{figure}[h]
    \begin{minipage}[c]{2.15cm}
        \centering
        \includegraphics[height=2cm,width=2cm]{cydQRcode.jpeg}
    \end{minipage}
    \hfill
    \begin{minipage}[c]{2.5cm}
        \centering
        \includegraphics[height=3.5cm,width=2.5cm]{cyd.jpg}
    \end{minipage}
\end{figure}
\vspace{-28mm}
% ========================================================================
\name{\Large\textbf{\fangsong 陈胤达}}
% ========================================================================
\basicInfo{\phone{(+86)13058626611}\quad\email{cyd0806@mail.ustc.edu.cn}
            }

 \centering{\faChain \url{https://ydchen0806.github.io/}} \\
  \centering{\faChain \url{https://scholar.google.com/citations?user=hCvlj5cAAAAJ&hl=en&oi=ao}}
% ========================================================================
\vspace{.5em}

% ========================================================================
% ========================================================================
\section{\makebox[.75em][c]{\faGraduationCap} \textbf{\fangsong 教育背景}}
% ========================================================================
\datedsubsection{\textbf{中国科学技术大学 \& 上海AI Lab}\ {\faMapMarker \small\textbf{\fangsong 合肥 \& 上海}}}{}
% \datedsubsection{\textbf{中国科学技术大学}\ {\faMapMarker \small\textbf{\fangsong 合肥}}}{}
\begin{flushleft}
\textbf{\fangsong 信息与通信工程，博士研究生} \hfill \textbf{2024.7 \textasciitilde\ 2027.6（预计）}
\end{flushleft}
\begin{itemize}
    \item \textbf{\kaishu 研究方向:}\hangindent=2em
    \kaishu 预测性多模态智能：从自监督表征与预测学习到具身世界模型和VLA；兼顾科学智能体、基础模型后训练、医学影像与连接组分析
    \item \textbf{\kaishu 导师与科研指导:}\hangindent=2em
    \kaishu \href{https://scholar.google.com/citations?user=5bInRDEAAAAJ&hl=en}{吴枫}（中国工程院院士，IEEE Fellow），\href{https://scholar.google.com/citations?user=Snl0HPEAAAAJ&hl=en&oi=ao}{熊志伟}；联培导师（2023）：\href{https://scholar.google.com/citations?user=qpBtpGsAAAAJ&hl=en}{汤晓鸥}（IEEE Fellow，上海AI Lab）；连接组研究指导：张若冰（复旦大学）
    \item \textbf{\kaishu 核心课程:}\hangindent=2em
    \kaishu 算法设计与分析，统计学习，深度学习，强化学习
    \item \textbf{\kaishu 荣誉奖励:}\hangindent=2em
    \kaishu 国家自然科学基金博士生项目负责人（2024年获批，2025--2027年执行）；\href{https://ydchen0806.github.io/docs/cast_ccf_2026_selection_proof.png}{入选中国科协博士生青托（2026，CCF培育）}
\end{itemize}

\datedsubsection{\textbf{中国科学技术大学}\ {\faMapMarker \small\textbf{\fangsong 合肥}}}{}
\begin{flushleft}
\textbf{\fangsong 计算机技术，硕士} \hfill \textbf{2022.9 \textasciitilde\ 2024.7}
\end{flushleft}
\begin{itemize}
    \item \textbf{\kaishu 荣誉奖励:}\hangindent=2em
    \kaishu 研究生国家奖学金（2022年）
\end{itemize}

\datedsubsection{\textbf{厦门大学}\ {\faMapMarker \small\textbf{\fangsong 厦门}}}{}
\begin{flushleft}
\textbf{\fangsong 环境生态工程\&经济学双学位，本科} \hfill \textbf{2018.9 \textasciitilde\ 2022.7}
\end{flushleft}
\begin{itemize}
    \item \textbf{\kaishu 综合排名: 1/31}
    % \textbf{\kaishu 专业排名: 16/31} \qquad
    \item \textbf{\kaishu 荣誉奖励:}\hangindent=2em
    \kaishu 厦门大学学术之星（2021年），CDA一级认证（2022年），Kaggle Expert
    \item \textbf{\kaishu 导师:}\hangindent=2em
    \kaishu \href{https://scholar.google.com/citations?user=l1GMXf4AAAAJ&hl=en&oi=ao}{张原野}
\end{itemize}
% ========================================================================
% ========================================================================
% ========================================================================
\section{\makebox[.75em][c]{\faGroup} \textbf{\fangsong 业界实习经历}}
% ========================================================================
\datedsubsection{\textbf{九坤投资（Ubiquant）}\ {\faMapMarker \small\textbf{\fangsong 北京}}}{}
\begin{flushleft}
\textbf{\fangsong 量化研究与大语言模型，研究实习生} \hfill \textbf{2026.5 \textasciitilde \ 2026.10}
\end{flushleft}
\begin{itemize}
    \item \kaishu 从事\textbf{量化研究}与\textbf{通用大模型}研发。
    \item \kaishu 具备\textbf{千卡集群}训练经验，探索通用大模型在数据驱动决策与因子挖掘中的应用。
    \item \kaishu \textbf{\href{https://arxiv.org/abs/2610.11451}{SAIL}核心贡献者}：面向文献研究、科学编程与多步骤科研任务（2026.10技术报告）。
    \item \kaishu 参与\textbf{\href{https://arxiv.org/abs/2610.30949}{IQuest-Q1}}智能体基座模型研究并署名技术报告（2026.10），面向命令行交互、代码与多步骤工具使用。
\end{itemize}

\datedsubsection{\textbf{北京人形机器人创新中心（X-Humanoid）}\ {\faMapMarker \small\textbf{\fangsong 北京}}}{}
\begin{flushleft}
\textbf{\fangsong 具身智能世界模型算法团队，核心成员} \hfill \textbf{2025.12 \textasciitilde \ 至今}
\end{flushleft}
\begin{itemize}
    \item \kaishu 研发面向人形机器人的\textbf{统一具身世界模型}，在真机闭环中联通任务相关感知、未来状态预测与低层动作生成。
    \item \kaishu 作为\textbf{Pelican-Unified 1.0}统一VLM与动作建模方向的核心贡献者，参与基模训练及规划验证；作为\textbf{Pelican-VLA 0.5}贡献者，研究任务相关瓶颈表征与动作学习。
    \item \kaishu 具备\textbf{百卡H200}集群训练经验，支持大规模具身基础模型的高效训练与优化。
\end{itemize}

\datedsubsection{\textbf{腾讯（IEG）}\ {\faMapMarker \small\textbf{\fangsong 上海}}}{}
\begin{flushleft}
\textbf{\fangsong 青云计划实习生} \hfill \textbf{2025.8 \textasciitilde \ 2025.12}
\end{flushleft}
\begin{itemize}
    \item \kaishu 参与\textbf{游戏世界模型}与多模态场景理解，包含和平精英项目\href{https://main-vla.github.io/}{\textbf{MAIN-VLA}}，联通意图、游戏状态与动作建模。
    \item \kaishu 落地王者荣耀\textbf{灵宝}项目，负责\textbf{TTS}与\textbf{实时对战解说}相关能力研发。
    \item \kaishu 基于F5-TTS和IndexTTS2训练并优化灵宝语音生成模型。
    \item \kaishu 合作开展\textbf{无路由MoE-LoRA}研究，探索低秩专家更新的线性聚合与基模合并。
\end{itemize}

\datedsubsection{\textbf{中国人民解放军总医院（301医院）}\ {\faMapMarker \small\textbf{\fangsong 北京}}}{}
\begin{flushleft}
\textbf{\fangsong 数据压缩小组，研究实习生} \hfill \textbf{2023.9 \textasciitilde \ 2024.2}
\end{flushleft}
\begin{itemize}
    \item \kaishu 协同戴琼海院士团队进行高效数据压缩的研究。
    \item \kaishu 研究医学影像特定的压缩算法，面向CT、MRI等模态优化内容表示、参考图像利用与编码效率。
\end{itemize}

% ========================================================================
\section{\makebox[.75em][c]{\faUniversity} \textbf{\fangsong 学界实习经历}}
% ========================================================================
\datedsubsection{\textbf{帝国理工学院}\ {\faMapMarker \small\textbf{\fangsong 伦敦（remote）}}}{}
\begin{flushleft}
\textbf{\fangsong Data Science Institute，RA} \hfill \textbf{2022.11 \textasciitilde \ 2023.8}
\end{flushleft}
\begin{itemize}
    \item \kaishu 与Rossella Arcucci、Che Liu合作开展三维医学多模态学习，形成GTGM生成文本引导预训练与BIMCV-R / MedFinder图文检索成果。
    \item \kaishu 开发图像--文本对比学习与检索框架；参与构建含\textbf{8,069个CT体积}的\href{https://github.com/ydchen0806/MedFinder}{BIMCV-R}数据集。
\end{itemize}

\datedsubsection{\textbf{厦门大学WISER Club}\ {\faMapMarker \small\textbf{\fangsong 厦门}}}{}
\begin{flushleft}
\textbf{\fangsong 数据挖掘小组，学术社群成员} \hfill \textbf{2021.8 \textasciitilde \ 2022.7}
\end{flushleft}
\begin{itemize}
    \item \kaishu 负责数据挖掘类课程的设计与讨论，主讲聚类和Transformer小节。
    \item \kaishu 指导本科生完成机器学习项目，并组织校内竞赛活动。
\end{itemize}

\datedsubsection{\textbf{厦门大学王亚南经济研究院}\ {\faMapMarker \small\textbf{\fangsong 厦门}}}{}
\begin{flushleft}
\textbf{\fangsong 计量经济学，研究助理} \hfill \textbf{2020.8 \textasciitilde \ 2021.12}
\end{flushleft}
\begin{itemize}
    \item \kaishu 协助朱炯副教授完成国土经济统计，进行宅基地信息的视觉特征提取。
    \item \kaishu 开发卫星影像分析工具，通过视觉特征分析土地利用变化。
\end{itemize}
% ========================================================================
% ========================================================================
\section{\makebox[.75em][c]{\faFlask} \textbf{\fangsong 科研项目}}
\datedsubsection{\textbf{\textit{具身智能与世界模型：Pelican 系列}}}{\textbf{2026.5 \textasciitilde\ 至今}}
\role {统一具身基础模型}{Pelican-Unified 1.0 核心贡献者；Pelican-VLA 0.5 贡献者}
\begin{itemize}
    \item \kaishu 在\href{https://arxiv.org/abs/2605.15153}{\textbf{Pelican-Unified 1.0}}中，作为\textbf{统一VLM与动作建模}方向核心贡献者，参与联通多模态理解、未来视频预测与低层动作生成，支持理解、推理、想象与动作的一体化建模。
    \item \kaishu 参与\href{https://arxiv.org/abs/2607.06655}{\textbf{Pelican-VLA 0.5: Attending Before Acting Benefits Generalization}}技术报告，围绕\textbf{任务相关瓶颈表征}、视觉语言条件与动作预测开展研究，并参与真实机器人闭环验证。
    \item \kaishu 方法重点：利用可学习\textbf{Bottleneck Token}汇聚任务相关视觉信息，研究未来状态建模与动作学习的联系，评估跨场景任务中的泛化表现。
    \item \kaishu 媒体报道：\href{https://www.peopleapp.com/column/30052155186-500007497131}{人民日报}，\href{https://www.beijing.gov.cn/fuwu/lqfw/gggs/202605/t20260519_4656998.html}{北京日报/北京市政府}，\href{https://mp.weixin.qq.com/s/z-OAqnl4IKHpaiT4bYBHIA}{机器之心}，\href{https://www.zhidx.com/p/558258.html}{智东西}
\end{itemize}

\datedsubsection{\textbf{\textit{国家自然科学基金青年学生基础研究项目（博士研究生）}}}{\textbf{2025.1 \textasciitilde \ 2027.12}}
\role {国家自然科学基金项目}{负责人\ {(2024年获批；执行期2025--2027年；经费\textcolor{red}{30万元})}}
\begin{enumerate}[label=\textbf{[\arabic*]}, leftmargin=2em, itemsep=3pt]
    \item \kaishu \textbf{神经科学表征与预训练}
    \begin{itemize}
        \item 面向电子显微镜与三维医学图像，研究视觉token构建、\textbf{自回归预训练}与多尺度表征一致性。
        \item 结合多模态神经影像数据，探索数据高效的\textbf{自监督与预测学习}方法。
    \end{itemize}
    \item \kaishu \textbf{精细神经元分割}
    \begin{itemize}
        \item 研究\textbf{多智能体掩码策略}、互补掩码一致性与跨域分割，提高表征迁移与分割鲁棒性。
    \end{itemize}
    \item \kaishu \textbf{连接组结构--功能分析}
    \begin{itemize}
        \item 分析成年果蝇蘑菇体连接组的神经递质侧化及其与记忆计算的联系。
        \item \textbf{SfN 2026 poster摘要已接收：}Connectome-scale chemical lateralization predicts asymmetric memory computation in the adult Drosophila mushroom body。
    \end{itemize}
\end{enumerate}

\datedsubsection{\textbf{\textit{表征与预测学习}}}{\textbf{2022.5 \textasciitilde \ 至今}}
\role{Conference\&Journal}{代表性论文（一作/共一）}
\begin{enumerate}[label=\textbf{[\arabic*]}, leftmargin=2em, itemsep=3pt]
    \item \kaishu \parbox[t]{0.88\linewidth}{\textbf{\href{https://openaccess.thecvf.com/content/ICCV2025/html/Chen_TokenUnify_Scaling_Up_Autoregressive_Pretraining_for_Neuron_Segmentation_ICCV_2025_paper.html}{TokenUnify: Scaling Up Autoregressive Pretraining for Neuron Segmentation}}, \textbf{ICCV 2025}} \hfill \textcolor{red}{\textbf{[共一]}}
    \begin{itemize}
        \item 提出结合\textbf{Mamba}的图像自回归预训练方法，面向长序列神经元图像学习可迁移的视觉表征。
        \item 设计\textbf{混合token预测策略}，结合局部结构与长程依赖，兼顾预训练效率与下游分割质量。
        \item 通过理论分析与实验研究模型规模扩展规律，验证预训练表征在电子显微镜神经元分割中的作用。
    \end{itemize}

    \item \kaishu \parbox[t]{0.88\linewidth}{\textbf{\href{https://www.ijcai.org/proceedings/2023/0068.pdf}{Self-supervised Neuron Segmentation with Multi-agent Reinforcement Learning}}, \textbf{IJCAI 2023}} \hfill \textcolor{red}{\textbf{[一作]}}
    \begin{itemize}
        \item 提出\textbf{dbMiM多智能体强化学习}框架，改进掩码自编码预训练，联合优化掩码率与掩码策略。
        \item 通过自适应掩码任务学习神经元图像的有效表征，并用于下游神经元分割。
    \end{itemize}

    \item \kaishu \parbox[t]{0.88\linewidth}{\textbf{\href{https://arxiv.org/abs/2308.09917}{Learning Multiscale Consistency for Self-Supervised Electron Microscopy Instance Segmentation}}, \textbf{ICASSP 2024}} \hfill \textcolor{red}{\textbf{[一作]}}
    \begin{itemize}
        \item 结合\textbf{多尺度特征对比学习}与特征重构，构建电子显微镜图像的自监督预训练策略。
        \item 设计\textbf{特征一致性损失}，缓解尺度变化引起的表征差异，并面向神经元实例分割进行验证。
    \end{itemize}

\end{enumerate}

\datedsubsection{\textbf{\textit{多模态与生成模型}}}{\textbf{2023.8 \textasciitilde \ 至今}}
\role{Conference\&Journal}{代表性论文（一作/共一）}
\begin{enumerate}[label=\textbf{[\arabic*]}, leftmargin=2em, itemsep=3pt]
    \item \kaishu \parbox[t]{0.88\linewidth}{\textbf{\href{https://doi.org/10.1109/TPAMI.2026.3744812}{Learned Image Coding with Generative Reference of Conditional Latents}}, \textbf{IEEE TPAMI, 2026 (accepted)}} \hfill \textcolor{red}{\textbf{[共一]}}
    \begin{itemize}
        \item 在CLC工作基础上，进一步研究\textbf{生成式参考图像}对条件潜变量编码与熵模型的帮助。
        \item 整合\textbf{本地字典、网络检索、图像生成}三类参考获取方式，研究多源参考与条件潜变量的联合建模。
        \item 结合参考图像合成、条件隐变量生成与鲁棒性分析，评估参考信息对率失真性能的作用。
    \end{itemize}

    \item \kaishu \parbox[t]{0.88\linewidth}{\textbf{\href{https://doi.org/10.1609/aaai.v39i12.33403}{Conditional Latent Coding with Learnable Synthesized Reference for Deep Image Compression}}, \textbf{AAAI 2025}} \hfill \textcolor{red}{\textbf{[共一]}}
    \begin{itemize}
        \item 构建图像相似度字典以检索参考图像，使用\textbf{可学习合成参考}改进熵模型的概率估计。
        \item 结合\textbf{条件潜变量编码}、参考信息合成与自适应熵建模，研究参考图像辅助的深度图像压缩。
    \end{itemize}

    \item \kaishu \parbox[t]{0.88\linewidth}{\textbf{\href{https://openaccess.thecvf.com/content/ICCV2025W/VLM3D/html/Chen_GTGM_Generative_Text-Guided_3D_Vision-Language_Pretraining_for_Medical_Image_Segmentation_ICCVW_2025_paper.html}{GTGM: Generative Text-Guided 3D Vision-Language Pretraining for Medical Image Segmentation}}, \textbf{ICCV 2025 Workshop (VLM3D)}} \hfill \textcolor{red}{\textbf{[共一]}}
    \begin{itemize}
        \item 利用大语言模型\textbf{生成图像描述}，在不依赖真实配对文本的条件下开展三维医学视觉语言预训练。
        \item 结合图文对比学习与生成文本引导，学习用于医学图像分割的多模态表征；在10个数据集上进行下游有监督微调与评估。
    \end{itemize}

    \item \kaishu \parbox[t]{0.88\linewidth}{\textbf{\href{https://arxiv.org/abs/2403.15992}{BIMCV-R: A Landmark Dataset for 3D CT Text-Image Retrieval}}, \textbf{MICCAI 2024}} \hfill \textcolor{red}{\textbf{[一作]}}
    \begin{itemize}
        \item 参与构建含\textbf{8,069个CT体积}的三维医学图文检索数据集，为体积影像与临床报告之间的检索提供基准。
        \item 开发图文对齐与检索框架\href{https://github.com/ydchen0806/MedFinder}{\textbf{MedFinder}}，支持三维CT语义检索与相关信息搜索。
    \end{itemize}

    \item \kaishu \parbox[t]{0.88\linewidth}{\textbf{\href{https://arxiv.org/abs/2412.19080}{MaskFactory: Towards High-quality Synthetic Data Generation for Dichotomous Image Segmentation}}, \textbf{NeurIPS 2024}} \hfill \textcolor{red}{\textbf{[共一]}}
    \begin{itemize}
        \item 通过\textbf{刚性与非刚性形变}编辑掩码，再利用ControlNet生成对应的掩码--图像对。
        \item 结合形变驱动的掩码编辑、生成条件控制与合成数据质量评估，研究合成训练数据对二分图像分割的帮助。
    \end{itemize}

    \item \kaishu \parbox[t]{0.88\linewidth}{\textbf{\href{https://ieeexplore.ieee.org/document/11205280}{EMPOWER: Evolutionary Medical Prompt Optimization With Reinforcement Learning}}, \textbf{IEEE JBHI, 2025 (Early Access)}} \hfill \textcolor{red}{\textbf{[共一]}}
    \begin{itemize}
        \item 结合领域知识与\textbf{强化学习}开展医疗提示词的进化优化，使用医学注意力、多维度质量评估、进化搜索与语义验证。
        \item 论文实验中事实错误减少24.7\%，盲测中的临床医生偏好提升15.3\%，评估提示优化对医疗回答质量的作用。
    \end{itemize}

\end{enumerate}

\datedsubsection{\textbf{\textit{大模型工程}}}{\textbf{2023.9 \textasciitilde \ 至今}}
\role{Projects}{核心成员}
\begin{itemize}
    \item \kaishu \textbf{基础模型系统：}开展图像编码、通用大模型训练与量化研究工作流建设，具备\textbf{千卡集群}训练经验，使用DDP与DeepSpeed优化分布式训练。
    \item \kaishu \textbf{具身与生物医学模型：}参与真机闭环、Pelican-Unified基模训练及具身训练基础设施建设，具备\textbf{百卡H200}训练经验；开展医学图像与神经元模型预训练，使用\textbf{64卡A40}训练环境。
\end{itemize}


% ========================================================================
\section{\makebox[.75em][c]{\faTrophy} \textbf{\fangsong 荣誉奖项}}
% ========================================================================
\vspace{.5mm}
\begin{itemize}[leftmargin=0em, itemsep=3pt]
    \item \textbf{\fangsong 入选中国科协博士生青托（CCF培育）} \hfill \textbf{2026.10}
    \begin{itemize}[leftmargin=1.5em]
        \item \kaishu 培育学会：中国计算机学会（CCF）
    \end{itemize}

    \item \textbf{\fangsong 中国科学技术大学增华奖学金} \hfill \textbf{2026}

    \item \textbf{\fangsong 交叉学科建模竞赛（ICM）} \hfill \textbf{2024}
    \begin{itemize}[leftmargin=1.5em]
        \item \kaishu Outstanding Winner
    \end{itemize}

    \item \textbf{\fangsong 国家自然科学基金青年学生基础研究项目（博士研究生）} \hfill \textbf{2024.12}
    \begin{itemize}[leftmargin=1.5em]
        \item \kaishu 项目负责人；2024年获批，执行期2025--2027年，经费30万元
    \end{itemize}

    \item \textbf{\fangsong 研究生国家奖学金} \hfill \textbf{2022.12}
    \begin{itemize}[leftmargin=1.5em]
        \item \kaishu 中国科学技术大学研究生国家奖学金
    \end{itemize}

    \item \textbf{\fangsong 厦门大学学术之星} \hfill \textbf{2021.12}
    \begin{itemize}[leftmargin=1.5em]
        \item \kaishu 本科阶段科研荣誉
    \end{itemize}

    \item \textbf{\fangsong ``景润杯"数学竞赛数学专业类} \hfill \textbf{2021.10}
    \begin{itemize}[leftmargin=1.5em]
        \item \kaishu 厦门大学第一名
    \end{itemize}

    \item \textbf{\fangsong ``互联网+"大赛} \hfill \textbf{2021.08}
    \begin{itemize}[leftmargin=1.5em]
        \item \kaishu 福建省金奖
    \end{itemize}

    \item \textbf{\fangsong 全国大学生数学竞赛非数学专业组决赛} \hfill \textbf{2021.05}
    \begin{itemize}[leftmargin=1.5em]
        \item \kaishu 全国二等奖
    \end{itemize}

    \item \textbf{\fangsong ``挑战杯"全国大学生课外学术科技作品竞赛} \hfill \textbf{2021.05}
    \begin{itemize}[leftmargin=1.5em]
        \item \kaishu 福建省一等奖
    \end{itemize}

    \item \textbf{\fangsong 全国大学生数学竞赛非数学专业组} \hfill \textbf{2020.11}
    \begin{itemize}[leftmargin=1.5em]
        \item \kaishu 福建省第一名
    \end{itemize}
    \end{itemize}
% ========================================================================
\section{\makebox[.75em][c]{\faCogs} \textbf{\fangsong 专业技能}}
% ========================================================================
\noindent
\begin{minipage}[t]{.5\textwidth}
    \begin{itemize}
    \item \textbf{\kaishu 编程能力}: Python, MATLAB, \LaTeX, C, C++, Java
    \item \textbf{\kaishu 深度学习}: PyTorch, TensorFlow, DeepSpeed, DDP
    \end{itemize}
\end{minipage}
\begin{minipage}[t]{.5\textwidth}
    \begin{itemize}
    \item \textbf{\kaishu 英语能力}: TOEFL(110), GRE(328)
    \item \textbf{\kaishu 专业工具}: Git, Docker, CUDA, HPC
    \end{itemize}
\end{minipage}

% ========================================================================
\section{\makebox[.75em][c]{\faUsers} \textbf{\fangsong 学术服务}}
% ========================================================================
\begin{itemize}
    \item \textbf{\kaishu 期刊\&会议审稿人}: CVPR 2025, ICCV 2025, WACV 2026, NeurIPS 2024, ICML 2025, ICLR 2024, MICCAI 2025, ACM MM 2024, AISTATS 2024, IJCV, TIP
    \item \textbf{\kaishu 邀请报告}: Vision-Language Pretraining Using Generative Methods，京东AI团队，2023.03。
\end{itemize}

% ========================================================================
\end{document}
